% Options for packages loaded elsewhere
\PassOptionsToPackage{unicode}{hyperref}
\PassOptionsToPackage{hyphens}{url}
\documentclass[
  ignorenonframetext,
]{beamer}
\newif\ifbibliography
\usepackage{pgfpages}
\setbeamertemplate{caption}[numbered]
\setbeamertemplate{caption label separator}{: }
\setbeamercolor{caption name}{fg=normal text.fg}
\beamertemplatenavigationsymbolsempty
% remove section numbering
\setbeamertemplate{part page}{
  \centering
  \begin{beamercolorbox}[sep=16pt,center]{part title}
    \usebeamerfont{part title}\insertpart\par
  \end{beamercolorbox}
}
\setbeamertemplate{section page}{
  \centering
  \begin{beamercolorbox}[sep=12pt,center]{section title}
    \usebeamerfont{section title}\insertsection\par
  \end{beamercolorbox}
}
\setbeamertemplate{subsection page}{
  \centering
  \begin{beamercolorbox}[sep=8pt,center]{subsection title}
    \usebeamerfont{subsection title}\insertsubsection\par
  \end{beamercolorbox}
}
% Prevent slide breaks in the middle of a paragraph
\widowpenalties 1 10000
\raggedbottom
\AtBeginPart{
  \frame{\partpage}
}
\AtBeginSection{
  \ifbibliography
  \else
    \frame{\sectionpage}
  \fi
}
\AtBeginSubsection{
  \frame{\subsectionpage}
}
\usepackage{iftex}
\ifPDFTeX
  \usepackage[T1]{fontenc}
  \usepackage[utf8]{inputenc}
  \usepackage{textcomp} % provide euro and other symbols
\else % if luatex or xetex
  \usepackage{unicode-math} % this also loads fontspec
  \defaultfontfeatures{Scale=MatchLowercase}
  \defaultfontfeatures[\rmfamily]{Ligatures=TeX,Scale=1}
\fi
\usepackage{lmodern}
\ifPDFTeX\else
  % xetex/luatex font selection
\fi
% Use upquote if available, for straight quotes in verbatim environments
\IfFileExists{upquote.sty}{\usepackage{upquote}}{}
\IfFileExists{microtype.sty}{% use microtype if available
  \usepackage[]{microtype}
  \UseMicrotypeSet[protrusion]{basicmath} % disable protrusion for tt fonts
}{}
\makeatletter
\@ifundefined{KOMAClassName}{% if non-KOMA class
  \IfFileExists{parskip.sty}{%
    \usepackage{parskip}
  }{% else
    \setlength{\parindent}{0pt}
    \setlength{\parskip}{6pt plus 2pt minus 1pt}}
}{% if KOMA class
  \KOMAoptions{parskip=half}}
\makeatother
\usepackage{longtable,booktabs,array}
\usepackage{caption}
\captionsetup[table]{skip=6pt}
\newcounter{none} % for unnumbered tables
\usepackage{calc} % for calculating minipage widths
\usepackage{caption}
% Make caption package work with longtable
\makeatletter
\def\fnum@table{\tablename~\thetable}
\makeatother
\setlength{\emergencystretch}{3em} % prevent overfull lines
\providecommand{\tightlist}{%
  \setlength{\itemsep}{0pt}\setlength{\parskip}{0pt}}
% Manuscript macro/package fallbacks (newcommand -> providecommand).
\usepackage{amsmath}
\usepackage{amssymb}
\usepackage{booktabs}
% Contents listing: article.cls prefixes every \l@section entry with
% \addvspace{1.0em}, which across ten sections costs about ten lines and pushed
% the ~40-entry listing one line past page 2 — leaving a third sheet carrying a
% single "10 References" line. Tighten that inter-section gap for the duration
% of \tableofcontents only, inside a group, so body spacing is untouched and no
% entry is dropped (reducing \tocdepth would have hidden 29 real subsections).
  \providecommand{\addvspace}[1]{\vskip 2\p@}%
% Document metadata. config.yaml declares a subtitle and a keyword list, and
% the producer emits both as tokens, but the template's \hypersetup writes only
% pdftitle/pdfauthor/pdflang -- so `pdfinfo` reported no Subject and no
% Keywords. This block runs after the template's, so it wins.
%
% The two values below are WRITTEN by scripts/z_generate_manuscript_variables.py
% (sync_preamble_metadata), never typed. A double-brace token cannot be used here:
% preamble.md is substituted into output/manuscript/ like any other section, but
% the template's _manuscript_source.py then copies the raw file back over that
% copy, so an unresolved token would reach hyperref and land in the PDF's
% metadata verbatim.
% Bounded identifier splitting. The template defines
%   \protected\def\breaktt#1{\begingroup\ttfamily\seqsplit{#1}\endgroup}
% and \seqsplit permits a break after EVERY character with no continuation cue.
% In the 2026-09-05 PDF that rendered `model_family_manifest.json` as
% `model_family_manifest.js` + `on` and `src/manuscript_variables.py` as
% `(src/manusc` + `ript_variables.py` — the first is itself a valid filename, so
% a reader cannot tell a line break from a name. Keep seqsplit's guarantee that
% no code token overflows the measure, but forbid breaks inside the last
% \GNNttTail characters (an extension is never orphaned) and inside the first
% \GNNttHead (a one- or two-letter stub is never left behind).

% Slide layout defaults for warning-clean scientific decks.
\usepackage{etoolbox}
\IfFileExists{xurl.sty}{\usepackage{xurl}}{}
\IfFileExists{seqsplit.sty}{\usepackage{seqsplit}}{\newcommand{\seqsplit}[1]{#1}}
\protected\def\breakseq#1{\seqsplit{#1}}
\protected\def\breaktt#1{\begingroup\ttfamily\seqsplit{#1}\endgroup}
\setlength{\emergencystretch}{6em}
\tolerance=5000
\hbadness=10000
\hfuzz=1pt
\setlength{\tabcolsep}{2pt}
\AtBeginEnvironment{longtable}{\tiny\renewcommand{\arraystretch}{0.86}\setlength{\tabcolsep}{1pt}}
\AtBeginEnvironment{tabular}{\tiny\renewcommand{\arraystretch}{0.86}\setlength{\tabcolsep}{1pt}}
\AtBeginEnvironment{equation}{\tiny}
\AtBeginEnvironment{equation*}{\tiny}
\AtBeginEnvironment{align}{\tiny}
\AtBeginEnvironment{align*}{\tiny}
\AtBeginEnvironment{itemize}{\footnotesize}
\AtBeginEnvironment{enumerate}{\footnotesize}
\AtBeginEnvironment{description}{\footnotesize}
\setbeamerfont{caption}{size=\tiny}
\setbeamerfont{caption name}{size=\tiny}
\setbeamerfont{normal text}{size=\small}
\setbeamerfont{frametitle}{size=\small}
\setbeamerfont{section title}{size=\footnotesize}
\setbeamerfont{subsection title}{size=\footnotesize}
\setbeamertemplate{section page}{%
  \centering
  \begin{beamercolorbox}[sep=12pt,center,wd=\paperwidth]{section title}
    \parbox{0.86\paperwidth}{\centering\usebeamerfont{section title}\insertsection\par}
  \end{beamercolorbox}
}
\setbeamertemplate{subsection page}{%
  \centering
  \begin{beamercolorbox}[sep=8pt,center,wd=\paperwidth]{subsection title}
    \parbox{0.86\paperwidth}{\centering\usebeamerfont{subsection title}\insertsubsection\par}
  \end{beamercolorbox}
}
\setlength{\abovecaptionskip}{2pt}
\setlength{\belowcaptionskip}{0pt}

% Natbib and cross-reference fallbacks — slides don't load natbib
% or cleveref, but combined-PDF manuscript prose may emit these
% commands. The fallback renders citations as a bracketed key list
% and cross-references as detokenized labels so slides don't fail on
% undefined control sequences or raw underscores. \providecommand is
% a no-op if packages load later.
\providecommand{\citep}[1]{[#1]}
\providecommand{\citet}[1]{#1}
\providecommand{\citealp}[1]{#1}
\providecommand{\citeauthor}[1]{#1}
\providecommand{\citeyear}[1]{#1}
\providecommand{\cref}[1]{\texttt{\detokenize{#1}}}
\providecommand{\Cref}[1]{\texttt{\detokenize{#1}}}
\providecommand{\crefrange}[2]{\texttt{\detokenize{#1}}--\texttt{\detokenize{#2}}}
\providecommand{\Crefrange}[2]{\texttt{\detokenize{#1}}--\texttt{\detokenize{#2}}}
\providecommand{\paragraph}[1]{\textbf{#1}\ }

\newtheorem{proposition}{Proposition}
\newtheorem{hypothesis}{Hypothesis}
\newtheorem{remark}[theorem]{Remark}
\newtheorem{axiom}{Axiom}
\newtheorem{property}{Property}
\usepackage{bookmark}
\IfFileExists{xurl.sty}{\usepackage{xurl}}{} % add URL line breaks if available
\urlstyle{same}
% fallback for those not using the hyperref driver hyperxmp:
\makeatletter
\@ifundefined{xmpquote}{\newcommand{\xmpquote}[1]{#1}}{}
\makeatother
\hypersetup{
  hidelinks,
  pdfcreator={LaTeX via pandoc}}

\author{\texorpdfstring{}{}}
\date{}

\begin{document}

\section{Supplemental Source Surface}\label{sec:source_surface}

\begin{frame}[fragile,allowframebreaks]{Supplemental Source Surface}
This supplement records the top-level source surfaces a reader or
reviewer should inspect before turning any prose in the main manuscript
into a verifiable claim. Each surface below is authored material under
version control; the generated artifacts it produces are described
separately so that the boundary between hand-written source and
reproducible output stays explicit.

The \texttt{src/} tree is the executable core of
GeneralizedNotationNotation: A Text Language for Active Inference
Models. It is organized into 31 Python packages spanning 582 source
files and roughly 194646 lines of code. Alongside these packages sit the
25 top-level step modules (0--24) that implement the numbered pipeline
as thin orchestrators delegating into the packages, each module owning
one stage of the progression from a parsed GNN text model through
visualization, type checking, code export, and executable cognitive
\end{frame}

\begin{frame}[fragile,allowframebreaks]{Supplemental Source Surface}
simulation. Every claim the manuscript makes about pipeline behavior
should be traceable to one of these step modules rather than to
descriptive prose alone.

The \texttt{input/} tree holds the model corpora that exercise the
pipeline. Its \texttt{gnn\_files/} subtree contains 10 corpus
directories organized by model family, together totaling 29 example
files; \breaktt{input/multi\_agent\_models/} holds the
\texttt{multiagent} family's model outside that count, and
\breaktt{input/recursive\_models/} is a reserved directory that ships
documentation only. A \breaktt{model\_family\_manifest.json} enumerates
the 9 registered families and their target directories. This manifest is
the authoritative registry that downstream steps and the manuscript
variable producer read when they report family counts and
cross-framework coverage; it should be consulted directly rather than
inferred from directory listings.
\end{frame}

\begin{frame}[fragile,allowframebreaks]{Supplemental Source Surface}

The \texttt{scripts/} tree contains the thin orchestrators that gate and
reproduce the project. These include the acceptance scripts that confirm
the pipeline runs end to end, the reliability gates that enforce
determinism and coverage expectations, and the manuscript variable
producer (\breaktt{src/manuscript\_variables.py} driven from this layer)
that emits the double-brace \texttt{\{\{...\}\}} token values consumed
throughout the manuscript. Treating these scripts as the source of
reproduction commands keeps reported numbers bound to what the code
actually computes.

\end{frame}

\begin{frame}[fragile,allowframebreaks]{Supplemental Source Surface}
The \texttt{doc/} tree is the prose and reference surface, comprising
616 files of specification, tutorial, and design documentation for the
GNN language and its Active Inference grounding (Smékal and Friedman
2023). It is the place to verify that a manuscript statement about GNN
syntax or semantics matches the documented language rather than a
convenient paraphrase.

The \texttt{output/} tree collects per-step pipeline artifacts: the data
dumps, intermediate representations, validation reports, and the 105
figure artifacts committed under it. Everything here is disposable and
reproducible from the surfaces above, so it should be read as evidence
of a run rather than as authored source. Notably,
\breaktt{output/data/manuscript\_variables.json} is where the
manuscript's substituted token values are materialized.
\end{frame}

\begin{frame}[fragile,allowframebreaks]{Authored Source Versus Generated
Output}
\protect\phantomsection\label{authored-source-versus-generated-output}
\texttt{src/}, \texttt{scripts/}, \texttt{input/}, \texttt{doc/} and
\texttt{manuscript/} are authored and reviewed; everything under
\texttt{output/} is generated and disposable, including
\breaktt{output/manuscript/}, which holds the token-substituted copies of
the \texttt{manuscript/} sources rather than the sources themselves. The
commands that regenerate the output surfaces are listed in
sec.~6: the smoke run for the per-step
artifacts, \breaktt{scripts/z\_generate\_manuscript\_variables.py} for
the token map and the hydrated sections,
\breaktt{scripts.manuscript\_build\_figures} for the 6 manuscript
figures, and the template render stage for the PDF. Which values become
tokens is not a matter of taste: any quantity the producer can derive
from a source surface is a token,
\breaktt{scripts/check\_manuscript\_tokens.py} fails the build when a
\end{frame}

\begin{frame}[fragile,allowframebreaks]{Authored Source Versus Generated
Output}
section hard-codes one instead, and the resulting 82-entry map is
committed at \breaktt{output/data/manuscript\_variables.json} for audit.
External references are resolved from
\breaktt{manuscript/references.bib}, which the same gate cross-checks
against every citation key in the prose. This manuscript quotes no
private or unpublished material.
\end{frame}

\end{document}
